% Einheit 3 von 3 – Pythagoras in Figuren und Körpern (bank/pythagoras/e3.jsonl, zusammenbau v0.1)
%% TODO Zweigzeile Teil 1: was hier gelernt wird (Satz in Schülersprache aus dem Einheitstitel) – Einheit 3
\einheitenkopf[e3]{Einheit 3 von 3 · Pythagoras in Figuren und Körpern}
\zweigzeile{ab Kl. 8, je nach Buch bis Kl. 9 · P10 oft}
\setcounter{aufgabe}{20}

%% TODO Titel: Ich-kann-Satz fehlt in der Bank (Platzhalter: Kettenname) – Nr. 21
%% TODO Anweisung: ein Satz über der ersten Teilaufgabe fehlt in der Bank – Nr. 21
\begin{aufgabe}{Ganz oder halb?}
\begin{teile}
\teil Kegel mit Durchmesser $18$ cm – welche Länge gehört ins Teildreieck aus Höhe, Radius und Mantellinie? \\ \kreuz{$18$ cm} \\ \kreuz{$9$ cm}

\kegel{1.5}{3}{}{}{}
\end{teile}
\end{aufgabe}

%% TODO Titel: Ich-kann-Satz fehlt in der Bank (Platzhalter: Kettenname) – Nr. 22
%% TODO Anweisung: ein Satz über der ersten Teilaufgabe fehlt in der Bank – Nr. 22
\begin{aufgabe}{Figuren und Körper}
%% TODO \rechenplatz{n}: Schrittzahl des Grundfalls fehlt in der Bank (nur bei fester Schreibform, 2.3 b)
\begin{teile}
\teil Raute mit Diagonalen: ein rechtwinkliges Teildreieck nachfahren. Welche Seiten sind Katheten, welche ist die Hypotenuse?

\raute[diagonalen]{4}{2.5}
\teil Gleichschenkliges Dreieck: Grundseite $18$ cm, Schenkel $41$ cm – Höhe $h$?

\dreieck{(0,0)}{(2.02,0)}{(1.01,4.5)}{$41$ cm}{}{$18$ cm}{}{}{}
h² = \leerfeld , h = \leerfeld[cm]
\teil Gleichschenkliges Dreieck: Grundseite $24$ cm, Höhe $35$ cm – Schenkel? \leerfeld[cm]
\teil Gleichschenkliges Trapez: parallele Seiten $50$ cm und $32$ cm, Höhe $40$ cm – Schenkel?

\trapez[hoehe=h]{5}{3.2}{4}
\leerfeld[cm]
\teil Parallelogramm: schräge Seite $5{,}3$ m. Der Fußpunkt der Höhe liegt auf der Verlängerung der Grundseite, $2{,}8$ m hinter der Ecke – Höhe? h = \leerfeld[m]
\teil Dreieck ABC mit stumpfem Winkel bei B. Die Höhe von C trifft die Verlängerung von AB im Punkt D. $CD = 3{,}6$ cm, $BC = 3{,}9$ cm – wie lang ist BD? BD = \leerfeld[cm]
\teil Rechteck $1{,}8$ m mal $8{,}0$ m – Diagonale? \leerfeld[m]
\teil $P(1 | 1)$ und $Q(4 | 5)$ – Länge der Strecke PQ? PQ = \leerfeld
\teil $E(-4 | -3)$ und $F(5 | 9)$ – Länge der Strecke EF? EF = \leerfeld
\teil Quadratische Pyramide: Grundkante $18$ cm, Höhe $40$ cm – Seitenhöhe $h_s$?

\pyramide{4}{4}{8.89}{$18$ cm}{}{$40$ cm}
$h_s$ = \leerfeld[cm]
\teil Kegel: Durchmesser $66$ cm, Höhe $56$ cm – Mantellinie $s$?

\kegel{1.77}{3}{}{$56$ cm}{$s$}
s = \leerfeld[cm]
\teil Kegel: Mantellinie $82$ cm, Radius $18$ cm – Höhe $h$?

\kegel{0.68}{3}{$18$ cm}{$h$}{$82$ cm}
h = \leerfeld[cm]
\teil Beschreibe ohne Zahlen, wie man die Seitenhöhe einer quadratischen Pyramide aus Grundkante und Körperhöhe berechnet.
\teil Ein Turm besteht aus einem Zylinder ($18$ m hoch) und einem Kegeldach mit Radius $2{,}0$ m und Mantellinie $5{,}2$ m. Wie hoch ist der Turm?

\zylinder{1}{3}{}{$18$ m}
\leerfeld[m]
\teil Quader $2$ dm mal $3$ dm mal $6$ dm – Raumdiagonale?

\quader{1.33}{2}{4}{}{}{}
\leerfeld[dm]
\teil Ein Pavillon hat ein kegelförmiges Zeltdach mit dem Durchmesser $7{,}2$ m. Die Höhe des Dachs ist bekannt, steht hier aber nicht. Beschreibe, wie man die Länge einer Zeltstange vom Dachrand bis zur Spitze berechnet \hfill (P10 2018 OS).

\kegel{1.8}{1.5}{}{$h$}{$s$}
\teil Ein Getreidesilo besteht aus einem Zylinder mit der Höhe $19{,}5$ m und einem Kegeldach mit Radius $3{,}9$ m und Mantellinie $6{,}2$ m. Berechne die Gesamthöhe des Silos auf eine Dezimale \hfill (P10 2026 FOR).

\zylinder{1}{3}{}{$19{,}5$ m}
\leerfeld[m]
\teil Ein zylindrisches Glas ist innen $11$ cm hoch und hat den Radius $2{,}8$ cm. Ein Strohhalm steht schräg von der Bodenkante zur gegenüberliegenden oberen Kante und ragt $4$ cm heraus. Berechne die Länge des Strohhalms auf eine Dezimale \hfill (P10 2022 OS).

\zylinder{0.76}{3}{$2{,}8$ cm}{$11$ cm}
\leerfeld[cm]
\teil Die Punkte R und S sind im Koordinatensystem eingezeichnet. Berechne die Länge der Strecke RS auf eine Dezimale \hfill (P10 2019 OS).

\begin{ksys}[xmin=-4,xmax=5,ymin=-4,ymax=4,ablesen] \punkt{-2}{3}{R} \punkt{4}{-1}{S} \end{ksys}
RS ≈ \leerfeld
\end{teile}
\end{aufgabe}

%% TODO Titel: Ich-kann-Satz fehlt in der Bank (Platzhalter: Kettenname) – Nr. 23
%% TODO Anweisung: ein Satz über der ersten Teilaufgabe fehlt in der Bank – Nr. 23
\begin{aufgabe}{Rampe, Leiter, Seil an der Wand mit Skizze}
\begin{teile}
\teil Ein Seil ist an einer Wand in $5{,}6$ m Höhe befestigt und $3{,}3$ m vor der Wand im Boden verankert. Wie lang ist das Seil?

\dreieckrw{2.95}{5}{$3{,}3$ m}{$5{,}6$ m}{?}
\leerfeld[m]
\end{teile}
\end{aufgabe}

%% TODO Titel: Ich-kann-Satz fehlt in der Bank (Platzhalter: Kettenname) – Nr. 24
%% TODO Anweisung: ein Satz über der ersten Teilaufgabe fehlt in der Bank – Nr. 24
\begin{aufgabe}{Figuren und Körper – Fehler finden, Begründen, Anwendung, Darstellungswechsel}
\begin{teile}
\teil Gleichschenkliges Dreieck: Grundseite $1{,}4$ m, Schenkel $2{,}5$ m. Nils rechnet: \rechnung{h^2 &= 2{,}5^2 - 1{,}4^2 \\ h^2 &= 4{,}29 \\ h &\approx 2{,}1 \text{ m}} Wo ist der Fehler?
\teil Warum rechnet man im gleichschenkligen Dreieck mit der halben Grundseite?
\teil Ein Regenschirm ist $62$ cm lang und lässt sich nicht kürzer machen. Passt er flach in einen Koffer, der innen $55$ cm lang und $35$ cm breit ist?
\teil Zeichne $A(-2 | -1)$ und $B(4 | 3)$ ein, dazu das rechtwinklige Dreieck mit waagerechter und senkrechter Kathete. Beschrifte die Kathetenlängen.

\begin{ksys}[xmin=-3,xmax=5,ymin=-2,ymax=4] \end{ksys}
\end{teile}
\end{aufgabe}
